\begin{proof}[Proof of \cref{obj:thm:unrestricted-near-complete-arrival-envelope}]
\begin{enumerate}
\item Choose the universal constant \(\overline C>0\) supplied by
\cref{obj:thm:unrestricted-mixed-count-upper}. For every admissible law \(P\) and the specified sampling law, that result gives
\[
E_P\!\left[
\left(\widehat\tau^{\mathrm{mix}}_{n,d,q}(\mathbf O_n)-\tau(P)\right)^2
\right]
\le \mathfrak U_{n,d,q}
\le \overline C\,r_{n,d,q},
\]
and
\[
E_P\!\left[
\left(\widehat\tau^{\mathrm{mix}}_{n,d,q}(\mathbf O_n)-\tau(P)\right)^2
\right]
\le
E_P\!\left[
E_\omega\!\left[
\left(\widetilde\tau(\mathbf O_n;\omega)-\tau(P)\right)^2
\,\middle|\,\mathbf O_n
\right]\right].
\]
Its envelope is exactly the branchwise envelope in the statement, and its auxiliary output includes the zero output on the Poisson tail. We retain this same constant throughout.
% lean: CausalSmith.Stat.MarRareqLogfrontier.unrestricted_mixed_count_upper

\item We first record how a deterministic risk bound yields a minimax bound. Write
\[
\mathcal O_d:=\{0,\ldots,d-1\}\times\{0,1\}^{4}
\]
for the observed-record space. For a bounded estimator kernel \(\mathsf T\), define its risk by
\[
\mathcal R_P(\mathsf T)
:=
\int_{\mathcal O_d^n}\int_{[-1,1]}
(h-\tau(P))^2\,\mathsf T(\mathbf o,dh)\,
\mathcal L_P(\mathbf O_n)(d\mathbf o).
\]
These risks are nonnegative.

Let \(f_{\mathrm{det}}:\mathcal O_d^n\to[-1,1]\) be measurable, and let \(b_{\mathrm{up}}\ge0\). Suppose
\[
E_P\!\left[
\left(f_{\mathrm{det}}(\mathbf O_n)-\tau(P)\right)^2
\right]\le b_{\mathrm{up}}
\qquad
\text{for every }P\in\mathcal M^{+}_{n,d,q}.
\]
The parameter-independent kernel
\[
\mathsf T_{\mathrm{det}}(\mathbf o,\cdot)
:=\operatorname{Dirac}\!\left(f_{\mathrm{det}}(\mathbf o)\right)
\]
belongs to the decision class in \cref{obj:def:unrestricted-minimax-risk}, and satisfies
\[
\mathcal R_P(\mathsf T_{\mathrm{det}})
=
E_P\!\left[
\left(f_{\mathrm{det}}(\mathbf O_n)-\tau(P)\right)^2
\right].
\]
Consequently,
\[
\mathfrak R^{+}_{n,d,q}
\le
\sup_{P\in\mathcal M^{+}_{n,d,q}}
E_P\!\left[
\left(f_{\mathrm{det}}(\mathbf O_n)-\tau(P)\right)^2
\right]
\le b_{\mathrm{up}}.
\]
For a nonempty model class, the last inequality follows by taking the supremum of the pointwise bounds. For an empty class, the real supremum is defined to be zero, and the inequality follows from \(b_{\mathrm{up}}\ge0\).
% lean: CausalSmith.Stat.MarRareqLogfrontier.unrestricted_deterministic_upper

\item Fix \(n,d\ge1\) and \(0<q\le1\). Binary potential outcomes and
\cref{obj:def:ate-functional} give
\[
\tau(P)=P(Y(1)=1)-P(Y(0)=1)\in[-1,1].
\]
Indeed, each probability belongs to \([0,1]\). The zero map is measurable and takes values in \([-1,1]\), and its risk under the canonical product probability law is
\[
E_P\!\left[(0-\tau(P))^2\right]=\tau(P)^2\le1.
\]
Applying the deterministic comparison just established yields
\[
\mathfrak R^{+}_{n,d,q}\le1.
\]
% lean: CausalSmith.Stat.MarRareqLogfrontier.ate_mem_unit_interval

\item The rate in \cref{obj:def:frontier-rate} satisfies
\[
r_{n,d,q}
=
\min\left\{
1,\,
\frac1{nq}
+
\left[\frac{d}{nq\log(e+nq)}\right]^2
\right\}\ge0,
\qquad
\overline C\,r_{n,d,q}\ge0.
\]
The estimator in \cref{obj:def:mixed-count-estimator} is measurable. Its auxiliary output lies in \([-1,1]\) for every sample and every randomization, so averaging gives
\[
-1
\le
E_\omega\!\left[
\widetilde\tau(\mathbf o;\omega)\mid\mathbf o
\right]
=
\widehat\tau^{\mathrm{mix}}_{n,d,q}(\mathbf o)
\le1.
\]
Under the canonical product law, the coordinate observations are measurable and independent, with
\[
\mathcal L_P(O_i)=\mathcal L_P(O_1),
\qquad i=1,\ldots,n,
\]
so they satisfy \cref{obj:ass:iid-sampling}. The first step therefore applies to every law in the model class. The deterministic comparison, with \(b_{\mathrm{up}}=\overline C\,r_{n,d,q}\), gives
\[
\mathfrak R^{+}_{n,d,q}\le\overline C\,r_{n,d,q}.
\]
% lean: CausalSmith.Stat.MarRareqLogfrontier.unrestricted_near_complete_arrival_envelope

\item We next establish the observed-contrast bound. Fix
\(P\in\mathcal M^{+}_{n,d,q}\). For an observed record, write its coordinates as
\[
o=(X(o),A(o),S(o),R(o),(RY)(o)),
\]
and define
\[
\xi_{\mathrm{obs}}(o)
:=
2(2A(o)-1)\mathbf1\{R(o)=1,\ (RY)(o)=1\}.
\]
For a sample \(\mathbf o=(o_1,\ldots,o_n)\in\mathcal O_d^n\), define
\[
\overline\xi_n(\mathbf o)
:=
\frac1n\sum_{i=1}^n\xi_{\mathrm{obs}}(o_i).
\]
The canonical product law gives
\[
E_P[\overline\xi_n(\mathbf O_n)]
=E_P[\xi_{\mathrm{obs}}(O_1)],
\qquad
\operatorname{Var}_P(\overline\xi_n(\mathbf O_n))
=\frac1n\operatorname{Var}_P(\xi_{\mathrm{obs}}(O_1)).
\]
Since every observed coordinate is binary,
\[
-2\le\xi_{\mathrm{obs}}(o)\le2,
\qquad
\operatorname{Var}_P(\xi_{\mathrm{obs}}(O_1))
\le E_P[\xi_{\mathrm{obs}}(O_1)^2]\le4.
\]
Expanding the squared error around the mean now gives
\[
E_P\!\left[
\left(\overline\xi_n(\mathbf O_n)-\tau(P)\right)^2
\right]
=
\frac{\operatorname{Var}_P(\xi_{\mathrm{obs}}(O_1))}{n}
+
\left(E_P[\xi_{\mathrm{obs}}(O_1)]-\tau(P)\right)^2.
\]
By \cref{obj:def:complete-arrival-ht-estimator},
\[
\widehat\tau^{\mathrm{HT}}_n(\mathbf o)
=
\operatorname{clip}_{[-1,1]}(\overline\xi_n(\mathbf o)).
\]
Because \(\tau(P)\in[-1,1]\), considering whether
\(\overline\xi_n(\mathbf o)<-1\), belongs to \([-1,1]\), or exceeds \(1\), gives
\[
\left(
\widehat\tau^{\mathrm{HT}}_n(\mathbf o)-\tau(P)
\right)^2
\le
\left(\overline\xi_n(\mathbf o)-\tau(P)\right)^2.
\]
% lean: CausalSmith.Stat.MarRareqLogfrontier.observed_contrast_risk_le

\item To bound the bias, define the missing-outcome masses for both arms by
\[
\Delta_a(P)
:=
2P(A=a,R=0,Y=1),
\qquad a\in\{0,1\}.
\]
For every cell indexed by \(a,x,s\), the occupied-cell arrival condition implies
\[
q\,p_{axs}(P)
\le P(A=a,X=x,S=s,R=1).
\]
For an occupied cell this is the arrival-floor inequality in
\cref{obj:ass:occupied-cell-arrival}; for a null cell its left side is zero and its right side is nonnegative. Partitioning a cell according to \(R\) therefore gives
\[
\begin{aligned}
P(A=a,X=x,S=s,R=0)
&=
p_{axs}(P)-P(A=a,X=x,S=s,R=1)\\
&\le(1-q)p_{axs}(P).
\end{aligned}
\]
Balanced assignment gives \(P(A=1)=1/2\) and
\(P(A=0)=1-P(A=1)=1/2\). Thus, for each \(a\in\{0,1\}\),
\[
\sum_{x=0}^{d-1}\sum_{s=0}^{1}p_{axs}(P)=\frac12,
\]
and the finite cell partition yields
\[
\begin{aligned}
P(A=a,R=0)
&=
\sum_{x=0}^{d-1}\sum_{s=0}^{1}
P(A=a,X=x,S=s,R=0)\\
&\le
(1-q)\sum_{x=0}^{d-1}\sum_{s=0}^{1}p_{axs}(P)
=\frac{1-q}{2}.
\end{aligned}
\]
The event \(\{A=a,R=0,Y=1\}\) is contained in \(\{A=a,R=0\}\), so
\[
0\le\Delta_a(P)\le2P(A=a,R=0)\le1-q.
\]

Outcome consistency, randomized independence, and balanced assignment give
\[
\begin{aligned}
P(A=a,Y=1)
&=P(A=a,Y(a)=1)\\
&=P(A=a)P(Y(a)=1)
=\frac12P(Y(a)=1).
\end{aligned}
\]
Using the observed-record map and partitioning the outcome-one event according to arrival, we obtain
\[
\begin{aligned}
E_P[\xi_{\mathrm{obs}}(O_1)]
&=
2P(A=1,R=1,Y=1)-2P(A=0,R=1,Y=1)\\
&=
2P(A=1,Y=1)-2P(A=0,Y=1)
-\Delta_1(P)+\Delta_0(P)\\
&=\tau(P)-\Delta_1(P)+\Delta_0(P).
\end{aligned}
\]
The bounds on the two deficits imply
\[
-(1-q)
\le
E_P[\xi_{\mathrm{obs}}(O_1)]-\tau(P)
\le1-q,
\]
and hence, since \(1-q\ge0\),
\[
\left(E_P[\xi_{\mathrm{obs}}(O_1)]-\tau(P)\right)^2
\le(1-q)^2.
\]
Together with the preceding variance and clipping bounds, this gives
\[
\begin{aligned}
E_P\!\left[
\left(\widehat\tau^{\mathrm{HT}}_n(\mathbf O_n)-\tau(P)\right)^2
\right]
&\le
E_P\!\left[
\left(\overline\xi_n(\mathbf O_n)-\tau(P)\right)^2
\right]\\
&\le\frac4n+(1-q)^2.
\end{aligned}
\]
The clipped contrast is measurable on the finite sample space and takes values in \([-1,1]\). Since \(4/n+(1-q)^2\ge0\), the deterministic comparison gives both
\[
\mathfrak R^{+}_{n,d,q}\le\frac4n+(1-q)^2
\]
and
\[
\sup_{P\in\mathcal M^{+}_{n,d,q}}
E_P\!\left[
\left(\widehat\tau^{\mathrm{HT}}_n(\mathbf O_n)-\tau(P)\right)^2
\right]
\le\frac4n+(1-q)^2.
\]
% lean: CausalSmith.Stat.MarRareqLogfrontier.observed_contrast_risk_le

\item For the lower bound, first construct a complete-arrival pair. Define
\[
\delta_{\mathrm{gap}}:=\frac1{4\sqrt n},
\qquad
0<\delta_{\mathrm{gap}}\le\frac14,
\qquad
\delta_{\mathrm{gap}}^2=\frac1{16n}.
\]
For \(\iota\in\{0,1\}\), let \(P_{\mathrm{test},\iota}\) be the full-data law specified by
\[
\begin{gathered}
X=0,\qquad S(0)=S(1)=0,\qquad Y(0)=0,\qquad R=1,\\
A\sim\operatorname{Bernoulli}(1/2),\qquad
Y(1)\sim\operatorname{Bernoulli}(1/2+\iota\delta_{\mathrm{gap}}),
\qquad A\perp\!\!\!\perp Y(1),\\
S=S(A)=0,\qquad Y=Y(A)=A\,Y(1).
\end{gathered}
\]
The baseline value \(0\) is available for every \(d\ge1\), and the outcome parameters lie in \([1/2,3/4]\). Independence and the fair treatment coin verify randomized independence and balanced assignment. The displayed definitions verify consistency. Constant arrival verifies arrival MAR and gives arrival probability one in every occupied cell. Therefore
\[
P_{\mathrm{test},0},P_{\mathrm{test},1}
\in\mathcal M^{+}_{n,d,1},
\qquad
\tau(P_{\mathrm{test},\iota})
=\frac12+\iota\delta_{\mathrm{gap}}.
\]

Define the observed laws and their sample laws by
\[
Q_{\mathrm{test},\iota}
:=\mathcal L_{P_{\mathrm{test},\iota}}(O_1),
\qquad
Q_{\mathrm{test},\iota}^{(n)}
:=Q_{\mathrm{test},\iota}^{\otimes n}.
\]
Their common support consists of the three records
\[
\begin{aligned}
o_{\mathrm{test},\mathrm{control}}&:=(0,0,0,1,0),\\
o_{\mathrm{test},\mathrm{zero}}&:=(0,1,0,1,0),\\
o_{\mathrm{test},\mathrm{one}}&:=(0,1,0,1,1),
\end{aligned}
\qquad
\mathcal Z_{\mathrm{test}}
:=
\{o_{\mathrm{test},\mathrm{control}},
o_{\mathrm{test},\mathrm{zero}},
o_{\mathrm{test},\mathrm{one}}\}.
\]
The atom probabilities are
\[
\begin{aligned}
Q_{\mathrm{test},\iota}(\{o_{\mathrm{test},\mathrm{control}}\})&=\frac12,\\
Q_{\mathrm{test},\iota}(\{o_{\mathrm{test},\mathrm{zero}}\})&=\frac14-\frac{\iota\delta_{\mathrm{gap}}}{2},\\
Q_{\mathrm{test},\iota}(\{o_{\mathrm{test},\mathrm{one}}\})&=\frac14+\frac{\iota\delta_{\mathrm{gap}}}{2},
\end{aligned}
\]
with probability zero elsewhere. In particular,
\(Q_{\mathrm{test},1}\ll Q_{\mathrm{test},0}\).

For finite laws with a positive reference law on their common support, define
\[
\chi^2(Q_{\mathrm{test},1}\Vert Q_{\mathrm{test},0})
:=
\sum_{o\in\mathcal Z_{\mathrm{test}}}
\frac{
\left(Q_{\mathrm{test},1}(\{o\})-Q_{\mathrm{test},0}(\{o\})\right)^2
}{
Q_{\mathrm{test},0}(\{o\})
}.
\]
Summing the squared probability ratios gives
\[
\begin{aligned}
1+\chi^2(Q_{\mathrm{test},1}\Vert Q_{\mathrm{test},0})
&=
\sum_{o\in\mathcal Z_{\mathrm{test}}}
\frac{Q_{\mathrm{test},1}(\{o\})^2}{Q_{\mathrm{test},0}(\{o\})}\\
&=
\frac{(1/2)^2}{1/2}
+
\frac{(1/4-\delta_{\mathrm{gap}}/2)^2}{1/4}
+
\frac{(1/4+\delta_{\mathrm{gap}}/2)^2}{1/4}\\
&=1+2\delta_{\mathrm{gap}}^2.
\end{aligned}
\]
Thus
\[
\chi^2(Q_{\mathrm{test},1}\Vert Q_{\mathrm{test},0})
=2\delta_{\mathrm{gap}}^2=\frac1{8n}.
\]
% lean: CausalSmith.Stat.MarRareqLogfrontier.completeFamilyLaw_obs_chi

\item Define the sample likelihood ratio on \(\mathcal O_d^n\) by
\[
L_{\mathrm{test},n}(\mathbf o)
:=
\begin{cases}
\displaystyle\prod_{i=1}^{n}
\frac{Q_{\mathrm{test},1}(\{o_i\})}
{Q_{\mathrm{test},0}(\{o_i\})},
&\mathbf o\in\mathcal Z_{\mathrm{test}}^n,\\
0,&\mathbf o\notin\mathcal Z_{\mathrm{test}}^n.
\end{cases}
\]
The product laws satisfy
\[
E_{Q_{\mathrm{test},0}^{(n)}}[L_{\mathrm{test},n}]=1,
\]
and independence gives
\[
\begin{aligned}
1+\chi^2(Q_{\mathrm{test},1}^{(n)}\Vert Q_{\mathrm{test},0}^{(n)})
&=
E_{Q_{\mathrm{test},0}^{(n)}}[L_{\mathrm{test},n}^2]\\
&=
\left(
1+\chi^2(Q_{\mathrm{test},1}\Vert Q_{\mathrm{test},0})
\right)^n.
\end{aligned}
\]
Here the product divergence is defined by the same finite-sum formula on
\(\mathcal Z_{\mathrm{test}}^n\). Since \(1+t\le e^t\) for \(t\ge0\),
\[
\chi^2(Q_{\mathrm{test},1}^{(n)}\Vert Q_{\mathrm{test},0}^{(n)})
\le e^{1/8}-1<1.
\]
For the numerical inequality, monotonicity gives \(e^{1/8}\le e^{1/2}\), while
\[
(e^{1/2})^2=e<3<4,
\]
so \(e^{1/8}<2\).

Define total variation for these sample laws by
\[
\operatorname{TV}(Q_{\mathrm{test},1}^{(n)},Q_{\mathrm{test},0}^{(n)})
:=
\frac12
\sum_{\mathbf o\in\mathcal Z_{\mathrm{test}}^n}
\left|
Q_{\mathrm{test},1}^{(n)}(\{\mathbf o\})
-
Q_{\mathrm{test},0}^{(n)}(\{\mathbf o\})
\right|.
\]
Cauchy--Schwarz applied to the likelihood ratio yields
\[
\begin{aligned}
\operatorname{TV}(Q_{\mathrm{test},1}^{(n)},Q_{\mathrm{test},0}^{(n)})
&=
\frac12E_{Q_{\mathrm{test},0}^{(n)}}|L_{\mathrm{test},n}-1|\\
&\le
\frac12\sqrt{
E_{Q_{\mathrm{test},0}^{(n)}}[(L_{\mathrm{test},n}-1)^2]
}\\
&=
\frac12\sqrt{
\chi^2(Q_{\mathrm{test},1}^{(n)}\Vert Q_{\mathrm{test},0}^{(n)})
}
\le\frac12.
\end{aligned}
\]
Consequently, for every measurable \(\mathcal E\subseteq\mathcal O_d^n\),
\[
\begin{aligned}
Q_{\mathrm{test},1}^{(n)}(\mathcal E^c)
+Q_{\mathrm{test},0}^{(n)}(\mathcal E)
&=
1-\left(
Q_{\mathrm{test},1}^{(n)}(\mathcal E)
-Q_{\mathrm{test},0}^{(n)}(\mathcal E)
\right)\\
&\ge
1-\operatorname{TV}(Q_{\mathrm{test},1}^{(n)},Q_{\mathrm{test},0}^{(n)})
\ge\frac12.
\end{aligned}
\]
% lean: CausalSmith.Stat.MarRareqLogfrontier.complete_arrival_product_testing_half

\item Take any parameter-independent bounded kernel \(\mathsf T\) in the decision class, and define its barycenter by
\[
b_{\mathsf T}(\mathbf o)
:=
\int_{[-1,1]}h\,\mathsf T(\mathbf o,dh).
\]
By \cref{obj:lem:randomized-kernel-barycenter}, this is a measurable map into \([-1,1]\), and
\[
E_{P_{\mathrm{test},\iota}}\!\left[
\left(b_{\mathsf T}(\mathbf O_n)-\tau(P_{\mathrm{test},\iota})\right)^2
\right]
\le
\mathcal R_{P_{\mathrm{test},\iota}}(\mathsf T),
\qquad \iota\in\{0,1\}.
\]
Define the two large-error events and the nearest-target test by
\[
\mathcal E_{\mathrm{test},\iota}
:=
\left\{\mathbf o:
\frac{\delta_{\mathrm{gap}}}{2}
\le
\left|b_{\mathsf T}(\mathbf o)-\tau(P_{\mathrm{test},\iota})\right|
\right\},
\qquad \iota\in\{0,1\},
\]
and
\[
\mathcal A_{\mathrm{test}}
:=
\left\{\mathbf o:
\left|b_{\mathsf T}(\mathbf o)-\tau(P_{\mathrm{test},1})\right|
\le
\left|b_{\mathsf T}(\mathbf o)-\tau(P_{\mathrm{test},0})\right|
\right\}.
\]
The target separation and the triangle inequality give, for every sample,
\[
\delta_{\mathrm{gap}}
\le
\left|b_{\mathsf T}(\mathbf o)-\tau(P_{\mathrm{test},0})\right|
+
\left|b_{\mathsf T}(\mathbf o)-\tau(P_{\mathrm{test},1})\right|.
\]
On \(\mathcal A_{\mathrm{test}}\), the first error is at least
\(\delta_{\mathrm{gap}}/2\); on its complement, the second error is at least
\(\delta_{\mathrm{gap}}/2\). Hence
\[
\mathcal A_{\mathrm{test}}\subseteq\mathcal E_{\mathrm{test},0},
\qquad
\mathcal A_{\mathrm{test}}^c\subseteq\mathcal E_{\mathrm{test},1}.
\]
Applying the testing bound and then event monotonicity gives
\[
\frac12
\le
Q_{\mathrm{test},1}^{(n)}(\mathcal A_{\mathrm{test}}^c)
+
Q_{\mathrm{test},0}^{(n)}(\mathcal A_{\mathrm{test}})
\le
Q_{\mathrm{test},1}^{(n)}(\mathcal E_{\mathrm{test},1})
+
Q_{\mathrm{test},0}^{(n)}(\mathcal E_{\mathrm{test},0}),
\]
so
\[
\max_{\iota\in\{0,1\}}
Q_{\mathrm{test},\iota}^{(n)}(\mathcal E_{\mathrm{test},\iota})
\ge\frac14.
\]
Integrating squared error over each large-error event yields
\[
\left(\frac{\delta_{\mathrm{gap}}}{2}\right)^2
Q_{\mathrm{test},\iota}^{(n)}(\mathcal E_{\mathrm{test},\iota})
\le
E_{P_{\mathrm{test},\iota}}\!\left[
\left(b_{\mathsf T}(\mathbf O_n)-\tau(P_{\mathrm{test},\iota})\right)^2
\right].
\]
Whether the probability for \(\iota=0\) is smaller than that for \(\iota=1\), or conversely, multiplication by the nonnegative factor
\((\delta_{\mathrm{gap}}/2)^2\) preserves their ordering. Therefore
\[
\begin{aligned}
\frac{\delta_{\mathrm{gap}}^2}{16}
&\le
\left(\frac{\delta_{\mathrm{gap}}}{2}\right)^2
\max_{\iota\in\{0,1\}}
Q_{\mathrm{test},\iota}^{(n)}(\mathcal E_{\mathrm{test},\iota})\\
&=
\max_{\iota\in\{0,1\}}
\left\{
\left(\frac{\delta_{\mathrm{gap}}}{2}\right)^2
Q_{\mathrm{test},\iota}^{(n)}(\mathcal E_{\mathrm{test},\iota})
\right\}\\
&\le
\max_{\iota\in\{0,1\}}
E_{P_{\mathrm{test},\iota}}\!\left[
\left(b_{\mathsf T}(\mathbf O_n)-\tau(P_{\mathrm{test},\iota})\right)^2
\right]\\
&\le
\max_{\iota\in\{0,1\}}
\mathcal R_{P_{\mathrm{test},\iota}}(\mathsf T).
\end{aligned}
\]
For every bounded kernel and every full-data law, bounded actions and
\(\tau(P)\in[-1,1]\) also give
\[
0\le\mathcal R_P(\mathsf T)\le4.
\]
Thus the worst-case risks are finite. The decision class is nonempty because it contains the zero estimator, and the complete-arrival class contains the constructed pair. Taking the supremum over that class and then the infimum over kernels gives
\[
\mathfrak R^{+}_{n,d,1}
\ge
\frac{\delta_{\mathrm{gap}}^2}{16}
=
\frac1{256n}.
\]
% lean: CausalSmith.Stat.MarRareqLogfrontier.complete_arrival_two_point_minimax_floor

\item Finally, every complete-arrival law belongs to the model with arrival floor \(q\). Indeed, all other conditions are shared, and for each occupied cell of such a law,
\[
q\,p_{axs}(P)
\le p_{axs}(P)
\le P(A=a,X=x,S=s,R=1).
\]
Together with \(0<q\le1\), this proves
\[
\mathcal M^{+}_{n,d,1}\subseteq\mathcal M^{+}_{n,d,q}.
\]
For every kernel \(\mathsf T\), class inclusion gives
\[
\sup_{P\in\mathcal M^{+}_{n,d,1}}\mathcal R_P(\mathsf T)
\le
\sup_{P\in\mathcal M^{+}_{n,d,q}}\mathcal R_P(\mathsf T).
\]
The risks lie in \([0,4]\), and both classes contain the complete-arrival pair. Taking infima over the same nonempty decision class therefore yields
\[
\frac1{256n}
\le
\mathfrak R^{+}_{n,d,1}
\le
\mathfrak R^{+}_{n,d,q}.
\]
Combining this with the three upper bounds proves
\[
\frac1{256n}
\le
\mathfrak R^{+}_{n,d,q}
\le
\min\left\{
1,\,
\overline C\,r_{n,d,q},\,
\frac4n+(1-q)^2
\right\}.
\]
The first step supplies the mixed-count envelope and its auxiliary-risk comparison, and the observed-contrast calculation supplies the asserted uniform witness bound.
% lean: CausalSmith.Stat.MarRareqLogfrontier.unrestricted_near_complete_arrival_envelope
\end{enumerate}
\end{proof}
